\section{A weighted energy certificate and its boundary cost}
\label{sec:weighted-certificate}

For the smoothed Liouville sum \(B\) of \eqref{eq:gamma}, define
\[
 J(\sigma)=\int_1^\infty |B(X)|^2X^{-2\sigma}\frac{\dd X}{X},
 \qquad
 H(s)=\int_0^1 B(X)X^{-s}\frac{\dd X}{X}.
\]
The second integral is entire, by the exponential decay at zero.
The first is an extended nonnegative real number.

\begin{theorem}[Finite weighted energy excludes right-hand zeros]
\label{thm:weighted-certificate}
If \(J(\sigma)<\infty\) for some \(\sigma>1/2\), then
\(\zeta(s)\ne0\) on \(\R s>\sigma\). On that half-plane,
\[
 \left|\Gamma(s)\frac{\zeta(2s)}{\zeta(s)}-H(s)\right|^2
 \le\frac{J(\sigma)}{2(\R s-\sigma)}.
\]
Consequently RH is equivalent to \(J(\sigma)<\infty\) for every
\(\sigma>1/2\).
\end{theorem}
\begin{proof}
The integral \(T(s)=\int_1^\infty B(X)X^{-s}\dd X/X\)
converges absolutely on \(\R s>\sigma\), by Cauchy--Schwarz,
and gives the displayed inequality for \(T\). It is holomorphic
there. For example, the same argument with factors
\((\log X)^k\) bounds all differentiated integrals locally.
Initially on \(\R s>\max(1,\sigma)\), \eqref{eq:gamma} gives
\(\zeta(s)(T(s)+H(s))=\Gamma(s)\zeta(2s)\).
The identity extends meromorphically. Its right side is
holomorphic and nonzero on \(\R s>\sigma>1/2\).
A zero of \(\zeta\) there is therefore impossible.
If the energies are finite at every \(\sigma>1/2\), these
half-planes cover \(\R s>1/2\); the functional equation gives RH.
Conversely RH gives \(B(X)\ll_\varepsilon X^{1/2+\varepsilon}\)
by Proposition~\ref{prop:smoothed}. For a fixed \(\sigma>1/2\),
choose \(0<\varepsilon<\sigma-1/2\) and integrate its square.
\end{proof}

\begin{proposition}[A pole forces a quantitative energy cost]
\label{prop:weighted-pole-cost}
Suppose that \(\rho\) is a zero of \(\zeta\) of order \(m\) with
\(\beta=\R\rho>1/2\), and put
\[
 a_\rho=\frac{m!\,\Gamma(\rho)\zeta(2\rho)}{\zeta^{(m)}(\rho)}.
\]
Then
\[
 \liminf_{u\downarrow0}u^{2m-1}J(\beta+u/2)\ge |a_\rho|^2.
\]
The statement uses the extended value \(+\infty\) when an
energy diverges. Independently of RH,
\[
 \liminf_{u\downarrow0}u J(1/2+u/2)
 \ge\left|\frac{\Gamma(1/2)}{2\zeta(1/2)}\right|^2>0.
\]
\end{proposition}
\begin{proof}
At \(\rho\), the quotient in Theorem~\ref{thm:weighted-certificate}
has leading term \(a_\rho(s-\rho)^{-m}\), since
\(\Gamma(\rho)\zeta(2\rho)\ne0\). Subtracting \(H\)
preserves it. When \(J(\beta+u/2)\) is finite, evaluate the
inequality at \(s=\rho+u\), where
\(2(\R s-\beta-u/2)=u\). This gives
\(u^{2m-1}J(\beta+u/2)\ge
|u^m(\Gamma(s)\zeta(2s)/\zeta(s)-H(s))|^2\).
The right side tends to \(|a_\rho|^2\).
When the energy is infinite the inequality is automatic.
For the final assertion use the pole of \(\zeta(2s)\) at
\(s=1/2\), whose residue in the quotient is
\(\Gamma(1/2)/(2\zeta(1/2))\). The alternating zeta series
shows that \(\zeta(1/2)<0\).\footnote{This real pole is
present even when RH holds. It gives the unavoidable
boundary cost of the energy criterion.}
\end{proof}

These are explicit tests for a proposed energy estimate.
No upper bound for \(J(\sigma)\) in all \(\sigma>1/2\)
is established here. Proving that finiteness remains the
arithmetic step.

\begin{figure}[ht]
\centering
\begin{tikzpicture}[>=Stealth,every node/.style={font=\small}]
\node (a) at (0,0) {$J(\sigma)<\infty$};
\node (b) at (5.0,0) {$T(s)$ holomorphic};
\node (c) at (5.0,-1.0) {$\zeta(s)\ne0$ for $\Re s>\sigma$};
\draw[->] (a) -- node[above] {Cauchy--Schwarz} (b);
\draw[->] (b) -- (c);
\end{tikzpicture}
\caption{The proved implication from weighted energy to a zero-free
half-plane. The hypothesis in the left box remains an arithmetic
requirement for each $\sigma>1/2$.}
\label{fig:weighted-energy-certificate}
\end{figure}
